\section*{Chapter \ref{ch:g10:machine-integers} --- Integers in the Machine}

\begin{solution}{exo:g10:machine-integers:1}
(Python prints them with single quotes.) \code{"0b1101"}, \code{"0xff"}, \code{10}, \code{255} and \code{"0b100000000"} ($2^8 = 256$ needs nine bits).
\end{solution}

\begin{solution}{exo:g10:machine-integers:2}
From $0$ to $2^{12} - 1 = 4095$: $4096$ different values.
\end{solution}

\begin{solution}{exo:g10:machine-integers:3}
$-5$ is \code{11111011} (see the example of the chapter), $-1$ is \code{11111111}, $-128$ is \code{10000000}, $127$ is
\code{01111111}. Neither $128$ nor $-129$ can be written: the range on $8$ bits is $-128$ to $127$.
\end{solution}

\begin{solution}{exo:g10:machine-integers:4}
\code{1 << 10} is $1024 = 1 \times 2^{10}$; \code{100 >> 2} is $25$, the quotient of $100$ by $4$; \code{7 >> 3} is $0$,
the quotient of $7$ by $8$; \code{5 << 1} is $10 = 5 \times 2$; \code{0b1011 >> 1} is \code{0b101}, that is $5$, the
quotient of $11$ by $2$.
\end{solution}

\begin{solution}{exo:g10:machine-integers:5}
The function of the chapter's listing. For $n = 0$ it returns \code{"0"}; without the special case, the loop
\code{while n > 0} would not run and the function would return the empty string, which is not a way of writing a
number. The test: \code{int(to_base(n, b), b) == n} for many \code{n} and \code{b} (the chapter's tests try $500$).
\end{solution}

\begin{solution}{exo:g10:machine-integers:6}
At most $20 + 13 = 33$ bits, by the proposition on the bits of a product; the largest numbers on $20$ and $13$ bits
reach it, as the listing printed. At least $32$ bits: a $20$-bit number is at least $2^{19}$ and a $13$-bit number at
least $2^{12}$, so the product is at least $2^{31}$, which needs $32$ bits; $2^{19} \times 2^{12}$ is an example.
\end{solution}

\begin{solution}{exo:g10:machine-integers:7}
\code{int(bits, 2)} gives the leftmost bit the weight $+128$ instead of $-128$, that is, it returns the unsigned
reading, which is $256$ too large for a negative value: $-5 + 256 = 251$. A correct decoder subtracts $2^8$ when the
sign bit is $1$: \code{u = int(bits, 2)}, then \code{return u - 2**len(bits) if bits[0] == "1" else u}
(\code{exercises/decode_bug.py}).
\end{solution}

\begin{solution}{exo:g10:machine-integers:8}
The three functions of the method box (\code{exercises/bit_ops.py}): \code{(x >> i) & 1 == 1}, \code{x | (1 << i)} and
\code{x & ~(1 << i)}. Clearing after setting gives the same result as clearing directly, since the clear decides bit
\code{i} alone and setting changed no other bit; the chapter's tests check it on $300$ random pairs.
\end{solution}

\begin{solution}{exo:g10:machine-integers:9}
$100 + 50 = 150$: the pattern is \code{10010110}. Its unsigned value is $150$, which fits (no carry); its signed
value is $150 - 256 = -106$: the signed reading \omterm{def:g10:machine-integers:overflow}{overflows}. For $-100 + (-50)$, the patterns are $156$ and $206$, whose
sum $362$ leaves the pattern $362 - 256 = 106$, \code{01101010}: the unsigned addition carries, and the signed result
reads $106$ instead of $-150$, so both readings \omterm{def:g10:machine-integers:overflow}{overflow}.
\end{solution}

\begin{solution}{exo:g10:machine-integers:10}
\code{a = 1100}, \code{b = 1010}. After the first line \code{a = 0110}; after the second
\code{b = 0110 ^ 1010 = 1100}; after the third \code{a = 0110 ^ 1100 = 1010}. Now \code{a} is $10$ and \code{b} is
$12$: the values are swapped. With $a_0$, $b_0$ the starting values, the second line computes
$(a_0 \oplus b_0) \oplus b_0 = a_0 \oplus 0 = a_0$, and the third $(a_0 \oplus b_0) \oplus a_0 = b_0$, writing
$\oplus$ for \textsc{xor} (it can be computed in any order, as one checks bit by bit).
\end{solution}

\begin{solution}{exo:g10:machine-integers:11}
\code{(c >> 16) & 0xFF} is \code{0xF4} $= 244$, the red byte. \code{c >> 8} is \code{0xF4A2} $= 62{,}626$, the red and
green bytes together. Green is \code{(c >> 8) & 0xFF} $= 162$: \omterm{def:g10:machine-integers:shift}{shift} by one byte, then keep one byte.
\end{solution}

\begin{solution}{exo:g10:machine-integers:12}
Let $u$ be the unsigned reading of $x$ and $\bar u$ that of the flipped pattern. In each column one of the two has a
$1$, so $u + \bar u = 2^n - 1$ and $\bar u + 1 = 2^n - u$. If $x \ge 0$ then $u = x$ and $2^n - x$ is the pattern of
$-x$ (for $x = 0$ it is $2^n$, which keeps $n$ bits as $0$). If $x < 0$ then $u = x + 2^n$ and $2^n - u = -x$, which is
positive and in range as long as $x \ne -2^{n-1}$. For $x = -2^{n-1}$ (on $4$ bits, \code{1000}, that is $-8$),
flipping gives \code{0111} and adding $1$ gives \code{1000} again: $-(-8) = 8$ is out of range, and the negation
\omterm{def:g10:machine-integers:overflow}{overflows} back to $-8$.
\end{solution}

\begin{solution}{exo:g10:machine-integers:13}
Subtracting $1$ turns the rightmost $1$ bit of $x$ into $0$ and the $0$ bits to its right into $1$; the bits further
left do not change. The \textsc{and} with $x$ keeps those left bits and clears all the others, that is, it clears
exactly the rightmost $1$. So the second function turns its loop once per $1$ bit, while the first turns once per bit
of $x$ (\code{x.bit_length()} times). For $x = 2^{40}$: $41$ turns against $1$ (\code{exercises/popcount.py}).
\end{solution}

\begin{solution}{exo:g10:machine-integers:14}
If $x \ge 0 > y$, then $y \le x + y < x$, so the sum lies between two numbers in range and is in range. If both are
$\ge 0$, the true sum is at most $2(2^{n-1} - 1) = 2^n - 2$; when it is too large, it lies from $2^{n-1}$ to $2^n - 2$,
fits in $n$ unsigned bits, and its bit $n-1$ is $1$: the stored result reads as negative. If both are negative, the
true sum is at least $-2^n$; when it is too small, it lies from $-2^n$ to $-2^{n-1} - 1$, and the stored pattern, the
true sum plus $2^n$, lies from $0$ to $2^{n-1} - 1$: its sign bit is $0$. Conversely, when there is no \omterm{def:g10:machine-integers:overflow}{overflow} the
result has the sign of two operands of the same sign.
\end{solution}

\begin{solution}{exo:g10:machine-integers:15}
The file \code{exercises/add_n_bits.py} computes \code{total = a + b}, \texttt{result = total \% 2**n}, \code{carry = total >= 2**n},
and with \code{sign = 2**(n - 1)}, the \omterm{def:g10:machine-integers:overflow}{overflow} flag is true when \code{a & sign == b & sign} and
\code{result & sign != a & sign} (the proposition proved in the previous exercise). The tests compare it with
Python's arithmetic on $256 \times 37$ pairs of $8$-bit values. \code{add_n_bits(100, 50, 8)} returns
\code{(150, False, True)}.
\end{solution}

\begin{solution}{pb:g10:machine-integers:1}
\textbf{1.} $24 \times 3600 = 86{,}400$ seconds in a day; $365 \times 86{,}400 = 31{,}536{,}000$ in a year of $365$ days.

\textbf{2.} The leap years are 1972, 1976, 1980, 1984, 1988, 1992 and 1996: $30 \times 365 + 7 = 10{,}957$ days.

\textbf{3.} $10{,}957 \times 86{,}400 = 946{,}684{,}800$ seconds.

\textbf{4.} Python prints \code{946684800.0}.

\textbf{5.} $2^{31} - 1 = 2{,}147{,}483{,}647$.

\textbf{6.} \code{0x7FFFFFFF}: a $0$ sign bit followed by $31$ ones.

\textbf{7.} A year of $365.25$ days lasts $31{,}557{,}600$ seconds; $2{,}147{,}483{,}647 / 31{,}557{,}600 \approx 68.05$
years.

\textbf{8.} $1970 + 68 = 2038$.

\textbf{9.} 19~January 2038, 03:14:07 UTC.

\textbf{10.} $2{,}147{,}483{,}647 = 24{,}855 \times 86{,}400 + 11{,}647$: $24{,}855$ whole days and $11{,}647$ seconds.

\textbf{11.} $11{,}647 = 3 \times 3600 + 14 \times 60 + 7$: 03:14:07, as printed.

\textbf{12.} $2^{31} = 2{,}147{,}483{,}648$, one more than the largest signed $32$-bit integer.

\textbf{13.} The pattern written \code{0x80000000} in hexadecimal, a $1$ followed by $31$ zeros, whose signed
value is
$-2^{31} = -2{,}147{,}483{,}648$.

\textbf{14.} On the $4$-bit wheel, the clock counts $0, 1, \dots, 7$, then the next step lands on $-8$: it jumps from
the largest value to the smallest.

\textbf{15.} 13~December 1901, 20:45:52 UTC.

\textbf{16.} The register holds $-2^{31}$ instead of $2^{31}$: the clock is $2^{32} = 4{,}294{,}967{,}296$ seconds
behind the true time.

\textbf{17.} $2^{31} = 24{,}855 \times 86{,}400 + 11{,}648$, one second more than in question 10. Going back
$24{,}855$ days from 1~January 1970 reaches 14~December 1901, 00:00:00, and going back $11{,}648$ seconds more, that
is 3~h 14~min 8~s, reaches 13~December 1901, 20:45:52.

\textbf{18.} At $2^{32} - 1$ seconds: 7~February 2106, 06:28:15 UTC. An unsigned count cannot hold any negative
value, so it could not record a date before 1970.

\textbf{19.} About $292$ billion years.

\textbf{20.} A signed $32$-bit count of seconds since 1~January 1970 holds its last value, $2^{31} - 1$, at
19~January 2038, 03:14:07 UTC; one second later it wraps to $-2^{31}$, which reads as 13~December 1901, 20:45:52 UTC.
\end{solution}
